\backmatter
\chapter{Semisolutions}

\parbf{\ref{ex:n(n+1)/2}.}
Choose a point $p\in\Omega$ and a chart at $p$.

Assume that the main theorem holds for some dimension $d$.
Then the components $g_{ij}=g(\partial_i,\partial_j)$ of any Riemannian metric can be written as
\[
\langle \partial_i w, \partial_j w \rangle = g_{ij},
\eqlbl{eq:<ww>=g}
\]
for some smooth map $w\:\Omega \to \RR^d$.

There are $\tbinom{n+1}{2}=n\cdot (n+1)/2$ components $g_{ij}$, each of which is a smooth function on~$\Omega$.
The equation \ref{eq:<ww>=g} has $d$ unknowns; namely, the coordinate functions of~$w$.
This suggests that $d$ should be at least $\tbinom{n+1}{2}$.
The remaining proof is very standard, we reduce the problem to an algebraic system between the \emph{jet spaces} of $w$ and $g$ at the point $p$; the jet space corresponds to the space of Taylor polynomials of a given degree at~$p$.

Denote by $J^m(g)$ the degree-$m$ Taylor polynomial of $g$ at~$p$.
This polynomial has $\tbinom{m+n}{n}$ coefficients in an $n \cdot (n+1)/2$-dimensional space.
(Each coefficient corresponds to a certain partial derivative of $g$ of order at most $m$ at $p$.)
Such polynomials form an open subset of a $\tbinom{n+1}{2} \cdot \tbinom{m+n}{n}$-dimensional space.

Similarly, let $J^m(w)$ be the degree-$m$ Taylor polynomial of $w$ at $p$.
This polynomial has $\tbinom{m+n}{n}$ coefficients in a $d$-dimensional space;
hence these polynomials form a space of dimension $d \cdot \tbinom{m+n}{n}$.

The equations \ref{eq:<ww>=g} define a smooth map $J^{m+1}(w)\mapsto J^m(g)$.
In other words, any partial derivative of $g$ of degree at most $m$ can be expressed in terms of partial derivatives of $w$ of degree at most $m+1$.
Since the system \ref{eq:<ww>=g} admits a solution for any metric $g$, there exists a $J^{m+1}(w)$ mapping to any given $J^m(g)$.
By Sard’s lemma,
\[
d \cdot \tbinom{m+n+1}{n} \ge \tbinom{n+1}{2} \cdot  \tbinom{m+n}{n}.
\]
Since
$\tbinom{m+n+1}{n}/\tbinom{m+n}{n} \to 1$ as  $m \to \infty$, we conclude that $d \ge \tbinom{n+1}{2}$.

\parit{Remark.}
The proof is taken from \cite[Appendix I]{gromov-rokhlin} and it works for general system of partial differential equations.

Equation \ref{eq:<ww>=g} makes perfect sense for $C^1$-smooth maps $w$, but the above proof essentially uses that $w$ is $C^m$-smooth for sufficiently large~$m$.
Without this assumption, the statement does not hold.

Indeed, the Nash--Kuiper theorem \cite{nash-1954,kuiper-1955} states that equation \ref{eq:<ww>=g} admits a $C^1$-smooth solution $w$ if $d = n+1$.
This bound is optimal, since if $d=n$, then the metric $g$ must be flat, which is not the case in general.
Note that for $d = n+1$, the system \ref{eq:<ww>=g} is overdetermined, yet still admits a $C^1$-smooth solution.

It might happen that the optimal dimension for $C^2$-embeddings is smaller than $\tbinom{n+1}{2}$;
at least the our argument gives a worse bound in this case.
This question mentioned by Mikhael Gromov and Vladimir Rokhlin \cite[]{gromov-rokhlin}.\footnote{Counting degrees of freedom of curvature tesor one can get $d\ge \tfrac1{6}\cdot n\cdot (n+5)$.
Indeed the curvature nensor has $\tfrac1{12}\cdot n^2\cdot(n^2-1)$ idependent components, and each quadratic form has $n(n+1)$ idependent components.
Applying Gauss formula, we get that the codimension of the embedding has to be at least $\tfrac1{6}\cdot n\cdot (n+1)$.}




\parbf{\ref{ex:notinduced}.}
Consider square of a parallel 1-form $\alpha$ on $\mathbb{T}^2$ with kernel in irrational direction.

\parbf{\ref{ex:length-preserving}.} The only-if part is trivial.

Apply the length-preserving property to all arcs of curves in the coordinate directions to show that the equality
\[g_{ii}=\langle \partial_i w, \partial_i w\rangle\]
holds almost everywhere, and therefore everywhere.

Observe that
\[2\cdot g_{ij}=g(\partial_i+\partial_j,\partial_i+\partial_j)-g(\partial_i,\partial_i)-g(\partial_j,\partial_j).\]
Using this formula together with the argument above for curves in the direction $\partial_i$, $\partial_j$ and $\partial_i+\partial_j$, show that the equality
\[g_{ij}=\langle \partial_i w, \partial_j w\rangle\]
holds almost everywhere, and therefore everywhere.

\parbf{\ref{ex:R2->R}.} Assume the contrary;
let $w\:\RR^2\to \RR$ be a length-preserving map.

Note that $w$ is Lipschitz.
By Rademacher's theorem \cite{rademacher}, the differential $d_xw$ is defined for  almost all $x$.
Argue as in \ref{ex:length-preserving} to show that $|\partial_1w|=|\partial_2w|=1$, but $\partial_1w\cdot \partial_2w=0$ almost everywhere and arrive at a contradiction.

\parit{Remark.} This argument also shows that there is no length-preserving map from an $n$-dimensional Riemannian manifold to $\RR^{n-1}$.
On the other hand, a theorem of Mikhael Gromov \cite[Section~2.4.11]{gromov-1986} implies that \textit{any $n$-dimensional Riemannian manifold admits a length-preserving map into $\RR^n$}.
This is a close relative of the Nash--Kuiper theorem \cite{nash-1954,kuiper-1955}.
Another proof of Gromov's theorem is outlined in the previously mentioned lecture notes \cite{petrunin-yashinski}.

\parbf{\ref{ex:oplus}+\ref{ex:plus-oplus}.} Spell out the definitions.

\parbf{\ref{ex:f-ominus-phi}.}
Straightforward calculations.


\parbf{\ref{ex:torus-riemn}.}
By the definition of a submanifold, we can cover $\Omega$ by rectangular charts in which $\Omega$ is expressed as the coordinate subspace of the first $k=\dim \Omega$ coordinates.
Note that in each such chart, the metric on $\Omega$ can be written as the restriction of a Riemannian metric on the entire chart.

Extend this collection of charts to an atlas of $\TT^n$ such that the images of the remaining charts do not intersect $\Omega$.
Equip these new charts with constant Riemannian metrics.

Choose a partition of unity subordinate to the chart covering, and use it to patch together the local metrics into a smooth Riemannian metric on $\TT^n$.

\parbf{\ref{ex:short-embedding}.}
Let $\iota\:\Omega\to\RR^d$ be an embedding provided by Whitney’s theorem.
Think of $\RR^d$ as the affine subspace $\{x_{d+1} = 1\}$ in $\RR^{d+1}$.
The desired embedding can be found among maps of the form $x \mapsto \phi \cdot \iota(x)$, where $\phi\:\Omega \to \RR$ is a smooth positive function.

Try writing down the necessary conditions and convince yourself that such a function $\phi$ exists.

\parbf{\ref{ex:sum}.}
Let $\DD \subset \RR^n$ be the unit $n$-disc, and let $\DD^\circ$ denote its interior.
Construct a smooth map $r\:\RR^n \to \SSS^n$ such that $r|_{\DD^\circ}$ is a smooth embedding, and the set $\RR^n \setminus \DD^\circ$ is mapped to the south pole.

Since each $D_{\mathfrak{i}}$ is a smooth $n$-disc, there exists a diffeomorphism $\iota_{\mathfrak{i}}\:\DD \to D_{\mathfrak{i}}$.
Then the composition $r_{\mathfrak{i}} = r \circ \iota_{\mathfrak{i}}^{-1}$ satisfies the first two conditions.

Now, let us construct a metrics $h_{\mathfrak{i}}$ on $\SSS^n$ that meet the last condition.
Start by choosing sufficiently small metrics $\hat h_{\mathfrak{i}}$ on $\SSS^n$ such that
\[
g > \hat g = \sum_{\mathfrak{i}} r_{\mathfrak{i}}^* \hat h_{\mathfrak{i}},
\]
and let $g_1 = g - \hat g$.

Choose a partition of unity $\phi_{\mathfrak{i}}$ subordinate to the covering $\{D_{\mathfrak{i}}^\circ\}$;
that is, $\sum \phi_{\mathfrak{i}} = 1$, each $\phi_{\mathfrak{i}} \ge 0$, and the support of $\phi_{\mathfrak{i}}$ lies in $D_{\mathfrak{i}}^\circ$ for each~$\mathfrak{i}$.
Denote by $s_{\mathfrak{i}}$ the inverse of the restriction $r_{\mathfrak{i}}|_{D_{\mathfrak{i}}^\circ}$.

Observe that $s_{\mathfrak{i}}^*(\phi_{\mathfrak{i}} \cdot g_1)$ is a nonnegative metric on $\SSS^n$ and $\phi_{\mathfrak{i}} \cdot g_1\z=r_{\mathfrak{i}}^*s_{\mathfrak{i}}^*(\phi_{\mathfrak{i}} \cdot g_1)$ for each~$\mathfrak{i}$.
Therefore, the metric
\[
h_{\mathfrak{i}} = \hat h_{\mathfrak{i}} + s_{\mathfrak{i}}^*(\phi_{\mathfrak{i}} \cdot g_1)
\]
is Riemannian.
Verify that each $h_{\mathfrak{i}}$ satisfies the last condition.


\parbf{\ref{ex:free-free}.}
Show and use that a smooth map $f\:\Omega\to\RR^d$ is free at $p\in\Omega$ if there is no affine subspace $A$ of dimension less than $\tfrac{n\cdot (n+3)}2$ that has \emph{second-order contact} with $f$ at $p$;
that is,
\[\rho\circ f(x)=o(|p-x|^2),\]
where $\rho(z)$ denotes the distance from $z$ to $A$.

Alternatively, express the partial derivatives in one chart in terms of the partial derivatives in another chart and draw the conclusion.

\parbf{\ref{ex:free-dim}.}
The definition of a free map requires the linear independence of certain partial derivatives.
Count them.

\parbf{\ref{ex:free-sum}.}
Spell the definition.

\parbf{\ref{ex:free-twists}.}
This can be done by computing the partial derivatives.
Alternatively, one may use the observation about the order of contact with an affine subspace from \ref{ex:free-free}.

In the higher-dimensional case, one may take
\[\bigoplus_{i\le j}\Theta_{x_i+x_j},\]
where $x_i$ are coordinate functions.
The same argument should work.

\parbf{\ref{ex:free-composition}.} Use the definition of free map and \ref{ex:free-free}.

\parbf{\ref{ex:hoelder-norm}.} Apply the product rule for derivatives.

\parbf{\ref{ex:embedding-in-ball}.}
Check that all constructions in the proof of the main theorem can be carried out in an arbitrarily small neighborhood of the origin.

\begin{wrapfigure}{r}{32mm}
\vskip-2mm
\centering
\includegraphics{mppics/pic-20}
\end{wrapfigure}

Alternatively, construct an isometric embedding of $\RR^d$ into a small neighborhood of the origin in $\RR^{2 \cdot d}$ (take product of embeddings $\RR\hookrightarrow\RR^2$ shown on the picture)
and compose it with an isometric embedding $(\Omega, g) \hookrightarrow \RR^d$ provided by the main theorem.

\parbf{\ref{ex:Mv}.}
Note that the equation can be written as $L(w)=\hat h$, where $\hat h=h-L(v)$.
Then argue as in Nash’s lemma.
